\slajdid{10-s025}
\begin{frame}[label=10-s025]
\gusttekst\setlength{\abovedisplayskip}{3pt}\setlength{\belowdisplayskip}{3pt}\setlength{\abovedisplayshortskip}{2pt}\setlength{\belowdisplayshortskip}{2pt}Da bi odredili $I_{OL}$ smatraćemo da je na izlazu napon logičke nule. Tada je tranzistor u zasićenju. Na izlaznom čvoru važi
\[I_{C1}=I_{RC}+I_O\]
Povećanjem izlazne struje povećava se i struja kolektora tranzistora, a da bi bio stabilan napon logičke nule, potrebno je da tranzistor ostane u zasićenju, odnosno potrebno je da bude $\beta_F I_B>I_C$. Znači
\[\beta_F I_{B1}>I_{C1}=I_{RC}+I_O\]
odnosno
\[I_O<\beta_F I_{B1}-I_{RC}\]
Kako je po izlaznoj konturi
\[I_{RC}=\frac{V_{CC}-V_O}{R_C}\]
a tranzistor je u zasićenju
\[I_{RC}=\frac{V_{CC}-V_O}{R_C}=\frac{V_{CC}-V_{CES}}{R_C}\]
Znači
\[I_O<\beta_F I_{B1}-I_{RC}=\beta_F I_{B1}-\frac{V_{CC}-V_{CES}}{R_C}\]
\end{frame}
